
\documentclass{article}
\usepackage{colortbl}
\usepackage{makecell}
\usepackage{multirow}
\usepackage{supertabular}

\begin{document}

\newcounter{utterance}

\centering \large Interaction Transcript for game `adventuregame', experiment `new{-}words\_home{-}delivery\_easy\_basic\_easy', episode 6 with Qwen3{-}32B{-}t0.0.
\vspace{24pt}

{ \footnotesize  \setcounter{utterance}{1}
\setlength{\tabcolsep}{0pt}
\begin{supertabular}{c@{$\;$}|p{.15\linewidth}@{}p{.15\linewidth}p{.15\linewidth}p{.15\linewidth}p{.15\linewidth}p{.15\linewidth}}
    \# & \multicolumn{2}{c}{Player} && \multicolumn{2}{c}{Game Master} \\
    \hline

    \theutterance \stepcounter{utterance}  
    & & & \multicolumn{4}{p{0.6\linewidth}}{
        \cellcolor[rgb]{0.9,0.9,0.9}{
            \makecell[{{p{\linewidth}}}]{
                \texttt{\tiny{[P1$\langle$GM]}}
                \texttt{You are playing a text adventure game. I will describe what you can perceive in the game. You write the single action you want to take in the game starting with >. Only reply with an action.} \\
\texttt{For example:} \\
\texttt{> open cupboard} \\
\\ 
\texttt{In addition to common actions, you can tratel. To tratel is to physically place something somewhere.} \\
\\ 
\texttt{Your goal for this game is: Put the unpish on the table, the plate on the secti and the mop on the secti.} \\
\\ 
\texttt{Once you have achieved your goal, write "> done" to end the game.} \\
\\ 
\texttt{You are in a nolog now. There are a unpish, a wardrobe, a bed and a pillow. The unpish is on the floor. The wardrobe is closed. The pillow is on the bed. There is a passage to a hallway here.} \\
            }
        }
    }
     \\ \\

    \theutterance \stepcounter{utterance}  
    & \multicolumn{4}{p{0.6\linewidth}}{
        \cellcolor[rgb]{0.561,0.737,0.561}{
            \makecell[{{p{\linewidth}}}]{
                \texttt{\tiny{[P1$\rangle$GM]}}
                \texttt{<think>} \\
\texttt{Okay, let's see. My goal is to put the unpish on the table, the plate on the secti, and the mop on the secti. Wait, but in the current room description, there's no mention of a table, a plate, a mop, or a secti. The room has a nolog (maybe a typo for "room"?), a wardrobe, a bed, a pillow, and an unpish on the floor. There's a passage to a hallway.} \\
\\ 
\texttt{Hmm, so I need to find the table, secti, and the plate and mop. Since the current room doesn't have those items, I probably need to go to the hallway first. The passage leads to a hallway, so maybe the other items are there. Let me check the hallway next. So the first action should be to go through the passage into the hallway. Let me make sure I'm not missing anything here. The unpish is on the floor, but I need to put it on the table. So I need to pick it up first, but maybe after moving to the hallway and finding the table. Wait, but if I can't see the table here, I need to go elsewhere. So first, go to the hallway. Then, look for the table, secti, plate, and mop. Once I find them, I can start moving the items. So the first step is to go through the passage. Let me confirm the available actions.} \\
            }
        }
    }
    & & \\ \\

    \theutterance \stepcounter{utterance}  
    & & & \multicolumn{2}{p{0.3\linewidth}}{
        \cellcolor[rgb]{0.9,0.9,0.9}{
            \makecell[{{p{\linewidth}}}]{
                \texttt{\tiny{[GM$|$GM]}}
                \texttt{Round: 0} \\
            }
        }
    }
    & & \\ \\

    \theutterance \stepcounter{utterance}  
    & & & \multicolumn{2}{p{0.3\linewidth}}{
        \cellcolor[rgb]{0.9,0.9,0.9}{
            \makecell[{{p{\linewidth}}}]{
                \texttt{\tiny{[GM$|$GM]}}
                \texttt{command\_tag\_missing} \\
            }
        }
    }
    & & \\ \\

    \theutterance \stepcounter{utterance}  
    & & & \multicolumn{2}{p{0.3\linewidth}}{
        \cellcolor[rgb]{0.9,0.9,0.9}{
            \makecell[{{p{\linewidth}}}]{
                \texttt{\tiny{[GM$|$GM]}}
                {'goal_states_achieved': [], 'game_successfully_finished': False}
            }
        }
    }
    & & \\ \\

\end{supertabular}
}

\end{document}
